% !TeX root = ../sections/02_exercises.tex

\begin{exercise}
	r>1に対して、$\lim_{n\rightarrow\infty}r^n=\infty$であることを$\epsilon N$論法を用いて証明せよ。
\end{exercise}

\begin{background}
	この問題では，数列が正の無限大に発散することを示す。

	\begin{definition}[数列が正の無限大に発散すること]
		数列 $\{a_n\}$ が正の無限大に発散するとは，
		任意の実数 $M>0$ に対して，ある自然数 $N$ が存在して，
		すべての $n\geq N$ に対して
		\[
			a_n>M
		\]
		が成り立つことをいう。
		このとき，
		\[
			\lim_{n\to\infty}a_n=\infty
		\]
		と書く。
	\end{definition}

	\begin{remark}
		今回の数列は
		\[
			a_n=r^n
		\]
		である。したがって示すべきことは，任意の $M>0$ に対して，
		$n$ を十分大きくすれば
		\[
			r^n>M
		\]
		となることである。
	\end{remark}

	\begin{remark}
		前問のように有限の極限値を示す場合は
		\[
			|a_n-\alpha|<\epsilon
		\]
		を考えた。しかし，今回は極限値が $\infty$ であるため，
		ある値に近づくことではなく，任意に大きい数 $M$ を
		最終的に超えることを示す。
	\end{remark}

	\subsection{正の無限大への発散を示すための道具}

	\begin{theorem}[アルキメデスの性質]
		任意の実数 $A$ に対して，ある自然数 $N$ が存在して
		\[
			N>A
		\]
		が成り立つ。
	\end{theorem}

	\begin{lemma}
		$r>1$ のとき，
		\[
			\log r>0
		\]
		である。
	\end{lemma}

	\begin{remark}
		$\log r>0$ であるため，不等式の両辺を $\log r$ で割っても，
		不等号の向きは変わらない。
	\end{remark}
\end{background}

\begin{strategy}
	示したいことは
	\[
		\lim_{n\to\infty}r^n=\infty
	\]
	である。

	まず，正の無限大への発散の定義に戻る。
	この問題では，任意の $M>0$ に対して，ある自然数 $N$ が存在し，
	すべての $n\geq N$ に対して
	\[
		r^n>M
	\]
	が成り立つことを示せばよい。

	つまり，目標は
	\[
		r^n>M
	\]
	となるような自然数 $N$ を，$M$ に応じて見つけることである。

	この不等式を $n$ について解くために，両辺の対数をとる。
	$r>1$ なので $r^n>0$ であり，また $M>0$ であるから，
	対数をとることができる。すると
	\[
		r^n>M
	\]
	は
	\[
		n\log r>\log M
	\]
	となる。

	ここで $r>1$ より
	\[
		\log r>0
	\]
	である。したがって，両辺を $\log r$ で割っても
	不等号の向きは変わらず，
	\[
		n>\frac{\log M}{\log r}
	\]
	を得る。

	よって，アルキメデスの性質より，
	\[
		N>\frac{\log M}{\log r}
	\]
	を満たす自然数 $N$ を取ればよい。

	このように選べば，すべての $n\geq N$ に対して
	\[
		n>\frac{\log M}{\log r}
	\]
	となる。したがって，両辺に正の数 $\log r$ をかけて
	\[
		n\log r>\log M
	\]
	を得る。これは
	\[
		\log(r^n)>\log M
	\]
	を意味するので，
	\[
		r^n>M
	\]
	が従う。

	したがって，任意の $M>0$ に対して，十分大きな $n$ では
	$r^n>M$ が成り立つ。ゆえに
	\[
		\lim_{n\to\infty}r^n=\infty
	\]
	である。
\end{strategy}